% Generated draft: requires mathematical and source review.
\documentclass[11pt]{article}
\usepackage[margin=25mm]{geometry}
\usepackage{fontspec}
\setmainfont{Latin Modern Roman}
\usepackage{amsmath,amssymb,mathtools,unicode-math}
\setmathfont{Latin Modern Math}
\usepackage{enumitem,needspace,xurl,hyperref}
\setlength{\parindent}{0pt}
\setlength{\parskip}{5pt}
\setlength{\emergencystretch}{3em}
\newcommand{\N}{\mathbb{N}}
\newcommand{\Z}{\mathbb{Z}}
\newcommand{\Q}{\mathbb{Q}}
\newcommand{\R}{\mathbb{R}}
\newcommand{\C}{\mathbb{C}}
\newcommand{\F}{\mathbb{F}}
\newcommand{\A}{\mathbb{A}}
\newcommand{\PP}{\mathbb{P}}
\newcommand{\E}{\mathbb{E}}
\newcommand{\eps}{\varepsilon}
\newcommand{\dd}{\,\mathrm{d}}
\begin{document}
\section*{UnsolvedMath ID: 30002114}
\textbf{Source problem code:} catalogue entry
\subsection{Definitions and mathematical statement}
Let \(\Omega\subset\R^2\) be a nonempty bounded open set. In the standard star-shaped interpretation of ``starlike'', there exists \(x_0\in\Omega\) such that
\[
(1-t)x_0+tx\in\Omega\qquad\text{for every }x\in\Omega\text{ and }t\in[0,1].
\]
The supplied record does not specify additional boundary regularity or a stronger meaning of starlikeness.

Let \(H_0^1(\Omega)\) be the closure of \(C_c^\infty(\Omega)\) in the Sobolev space \(H^1(\Omega)\). Define the Dirichlet Laplacian through the quadratic form
\[
q_\Omega[u]=\int_\Omega |\nabla u|^2\,\dd x,
\qquad u\in H_0^1(\Omega).
\]
Write its eigenvalues, in nondecreasing order and repeated according to multiplicity, as
\[
0<\lambda_1(\Omega)\leq\lambda_2(\Omega)\leq\cdots.
\]
Equivalently, an eigenvalue \(\lambda\) admits a nonzero \(u\in H_0^1(\Omega)\) satisfying
\[
\int_\Omega\nabla u\cdot\overline{\nabla v}\,\dd x
=\lambda\int_\Omega u\overline v\,\dd x
\qquad\text{for every }v\in H_0^1(\Omega).
\]
This variational definition specifies the Dirichlet boundary condition without assuming a smooth boundary.

\textbf{Source question.} How does the shape of a bounded planar starlike domain \(\Omega\) affect its Dirichlet eigenvalues, and what geometrically sharp upper bounds depending on its shape hold?

\textbf{Unresolved formulation.} The record supplies a research question rather than a uniquely specified conjecture. It does not identify the geometric quantities allowed in a bound, a normalization or constraint on the domains, the eigenvalue indices or spectral expressions to be bounded, or the intended meaning of sharpness and its extremal domains. Consequently, no particular inequality, optimal constant, or equality characterization can be stated from this record alone. The definitions above clarify the spectral objects but do not resolve these missing specifications.
\subsection{Short English statement}
The question asks how the geometry of a bounded planar starlike domain controls the eigenvalues of its Dirichlet Laplacian, especially through sharp upper bounds. The supplied fragment leaves the geometric data, spectral target, domain constraints, and sharpness criterion unspecified, so it does not determine a single precise conjecture.
\subsection{Sources}
\begin{itemize}
\item Local UnsolvedMath registry, permanent ID \(30002114\), source code \texttt{catalogue entry}, title ``Sharp Shape-Dependent Bounds for Dirichlet Eigenvalues'', in the Oberwolfach Reports Open Problems collection. The supplied proposed year is 2012; no individual proposer is identified.
\item Supplied report citation: \emph{Geometric Aspects of Spectral Theory (2013)}, Oberwolfach Reports, pp.\ 2013--2076, DOI \texttt{10.4171/owr/2012/33}. Supplied source link: \url{https://doi.org/10.4171/owr/2012/33}.
\item Supplied literature source, labeled ``OWR report'': \url{https://ems.press/content/serial-article-files/46403}.
\item Literature assessment dated exactly \texttt{2026-08-21}: the supplied curation labels the problem open, characterizes the extraction as a broad fragment without a unique theorem or conjecture, and reports that the report describes the subject as a long standing difficult problem. This is supplied dated curation, not an independently verified literature assessment; time-sensitive status and cited claims require independent verification.
\item The supplied statement-verification field reports that the official report was checked and the complete statement retained byte-for-byte, with statement status ``exact''. This is a supplied provenance claim, not independent verification in this draft, and does not settle the formulation gaps described above.
\end{itemize}
\end{document}
